\documentclass[a4paper, 14pt]{extarticle}

% Поля
%--------------------------------------
\usepackage{geometry}
\geometry{a4paper,tmargin=2cm,bmargin=2cm,lmargin=3cm,rmargin=1cm}
%--------------------------------------

\usepackage{float}
%Russian-specific packages
%--------------------------------------
\usepackage[T2A]{fontenc}
\usepackage[utf8]{inputenc} 
\usepackage[english, main=russian]{babel}
%--------------------------------------

\usepackage{textcomp}
\usepackage{xcolor}

% Красная строка
%--------------------------------------
\usepackage{indentfirst}               
%--------------------------------------
%Листинги
%--------------------------------------
\usepackage{listings}
\lstset{
  basicstyle=\ttfamily\footnotesize,
  breakatwhitespace=false,
  breaklines=true,
  captionpos=t,
  inputencoding=utf8,
  extendedchars=true,
  literate={а}{а}1{б}{б}1{в}{в}1{г}{г}1{д}{д}1{е}{е}1{ё}{ё}1{ж}{ж}1{з}{з}1{и}{и}1{й}{й}1{к}{к}1{л}{л}1{м}{м}1{н}{н}1{о}{о}1{п}{п}1{р}{р}1{с}{с}1{т}{т}1{у}{у}1{ф}{ф}1{х}{х}1{ц}{ц}1{ч}{ч}1{ш}{ш}1{щ}{щ}1{ъ}{ъ}1{ы}{ы}1{ь}{ь}1{э}{э}1{ю}{ю}1{я}{я}1,
  frame=single,
  keepspaces=true,
  keywordstyle=\bf,
  numbers=none,
  xleftmargin=25pt,
  xrightmargin=25pt,
  showspaces=false,
  showstringspaces=false,
  showtabs=false,
  stepnumber=1,
  tabsize=2,
  title=\lstname
}
%--------------------------------------
\usepackage[most]{tcolorbox}
\tcbuselibrary{listings,breakable}
\newtcblisting{breakablelisting}{
  listing only,
  breakable,
  enhanced,
  frame hidden,
  colback=white,
  left=8pt,
  overlay={\draw[line width=0.4pt] (frame.north west) -- (frame.south west);
          \draw[line width=0.4pt] ([xshift=3pt]frame.north west) -- ([xshift=3pt]frame.south west);},
  listing options={
    frame=none,
    numbers=none,
    mathescape=false,
    texcl=false,
    inputencoding=utf8,
    extendedchars=true,
    literate={а}{а}1{б}{б}1{в}{в}1{г}{г}1{д}{д}1{е}{е}1{ё}{ё}1{ж}{ж}1{з}{з}1{и}{и}1{й}{й}1{к}{к}1{л}{л}1{м}{м}1{н}{н}1{о}{о}1{п}{п}1{р}{р}1{с}{с}1{т}{т}1{у}{у}1{ф}{ф}1{х}{х}1{ц}{ц}1{ч}{ч}1{ш}{ш}1{щ}{щ}1{ъ}{ъ}1{ы}{ы}1{ь}{ь}1{э}{э}1{ю}{ю}1{я}{я}1,
    literate={\$}{{\$}}1
  }
}


%Graphics
%--------------------------------------
\usepackage{graphicx}
\usepackage{wrapfig}
%--------------------------------------
\graphicspath{{./lab5/}{./}}
% Полуторный интервал
%--------------------------------------
\linespread{1.3}                    
%--------------------------------------

%Выравнивание и переносы
%--------------------------------------
% Избавляемся от переполнений
\sloppy
% Запрещаем разрыв страницы после первой строки абзаца
\clubpenalty=10000
% Запрещаем разрыв страницы после последней строки абзаца
\widowpenalty=10000
%--------------------------------------

%Списки
\usepackage{enumitem}

%Подписи
\usepackage{caption} 

%Гиперссылки
\usepackage{hyperref}

\hypersetup {
	unicode=true
}

%Рисунки
%--------------------------------------
\DeclareCaptionLabelSeparator*{emdash}{~--- }
\captionsetup[figure]{labelsep=emdash,font=onehalfspacing,position=bottom}
%--------------------------------------

\usepackage{tempora}

%--------------------------------------

%%% Математические пакеты %%%
%--------------------------------------
\usepackage{amsthm,amsfonts,amsmath,amssymb,amscd}  % Математические дополнения от AMS
\usepackage{mathtools}                              % Добавляет окружение multlined
\usepackage[perpage]{footmisc}
%--------------------------------------

%--------------------------------------
%			НАЧАЛО ДОКУМЕНТА
%--------------------------------------

\begin{document}

%--------------------------------------
%			ТИТУЛЬНЫЙ ЛИСТ
%--------------------------------------
\begin{titlepage}
\thispagestyle{empty}
\newpage


%Шапка титульного листа
%--------------------------------------
\vspace*{-60pt}
\hspace{-65pt}
\begin{minipage}{0.3\textwidth}
\hspace*{-20pt}\centering
\includegraphics[width=\textwidth]{emblem}
\end{minipage}
\begin{minipage}{0.67\textwidth}\small \textbf{
\vspace*{-0.7ex}
\hspace*{-6pt}\centerline{Министерство науки и высшего образования Российской Федерации}
\vspace*{-0.7ex}
\centerline{Федеральное государственное автономное образовательное учреждение }
\vspace*{-0.7ex}
\centerline{высшего образования}
\vspace*{-0.7ex}
\centerline{<<Московский государственный технический университет}
\vspace*{-0.7ex}
\centerline{имени Н.Э. Баумана}
\vspace*{-0.7ex}
\centerline{(национальный исследовательский университет)>>}
\vspace*{-0.7ex}
\centerline{(МГТУ им. Н.Э. Баумана)}}
\end{minipage}
%--------------------------------------

%Полосы
%--------------------------------------
\vspace{-25pt}
\hspace{-35pt}\rule{\textwidth}{2.3pt}

\vspace*{-20.3pt}
\hspace{-35pt}\rule{\textwidth}{0.4pt}
%--------------------------------------

\vspace{1.5ex}
\hspace{-35pt} \noindent \small ФАКУЛЬТЕТ\hspace{80pt} <<Информатика и системы управления>>

\vspace*{-16pt}
\hspace{47pt}\rule{0.83\textwidth}{0.4pt}

\vspace{0.5ex}
\hspace{-35pt} \noindent \small КАФЕДРА\hspace{50pt} <<Теоретическая информатика и компьютерные технологии>>

\vspace*{-16pt}
\hspace{30pt}\rule{0.866\textwidth}{0.4pt}
  
\vspace{11em}

\begin{center}
\Large {\bf Домашняя работа №1} \\ 
\large {\bf по курсу <<Моделирование>>} \\
\end{center}\normalsize

\vspace{8em}


\begin{flushright}
  {Студент группы ИУ9-82Б Кожин Н.С. \hspace*{15pt}\\ 
  \vspace{2ex}
  Преподаватель Домрачева А.Б.\hspace*{15pt}}
\end{flushright}

\bigskip

\vfill
 

\begin{center}
\textsl{Москва 2026}
\end{center}
\end{titlepage}
%--------------------------------------
%		КОНЕЦ ТИТУЛЬНОГО ЛИСТА
%--------------------------------------

\renewcommand{\ttdefault}{pcr}

\setlength{\tabcolsep}{3pt}
\newpage
\setcounter{page}{2}
\section{Постановка задачи}

Рассматриваются студенты первого курса, сдающие зачёт по информатике. Доступно три попытки с вероятностями успеха 0,6, 0,25 и 0,1 соответственно. Ежедневно экзаменуются не более 20 человек, между попытками~--- пауза минимум один день. Цель~--- оценить число успешно сдавших за 6 рабочих дней зачётной недели. В эксперименте рассматривается 90 студентов.

\section{Построение модели}

В качестве математической модели взята многоканальная система массового обслуживания. Обслуживаемыми единицами выступают студенты; каждая их попытка попасть на экзамен считается отдельной заявкой. Выбрана дисциплина с приоритетом первичным заявкам: те, кто приходит впервые, обслуживаются раньше пересдающих. Обоснование~--- у первичных сдающих выше вероятность успеха, поэтому разумно дать им преимущество в очереди и повысить суммарный результат.

Приоритет задаётся номером попытки $P_1 \in \{1, 2, 3\}$ по правилу $\text{priority} = P_1$: меньшему $P_1$ соответствует более высокий приоритет. Объём обслуживания~--- 20 параллельных каналов (мест). Один цикл обслуживания длится 1440 минут. Неудачная попытка приводит к задержке на сутки, после чего заявка возвращается в очередь с $P_1$, увеличенным на 1.

Жёсткое ограничение~--- не более 120 попыток за зачётную неделю (6 дней по 20 слотов). Заявки, получающие слот после истечения шестидневного интервала, считаются опоздавшими и не допускаются к экзамену. Структура: источник~--- генератор заявок; очередь с приоритетом; набор из 20 каналов. Связи: поступление в очередь $\to$ ожидание свободного канала $\to$ обслуживание $\to$ ветвление (выход при успехе, возврат в очередь при неудаче с учётом паузы).

Формула для ожидаемого числа сдавших при отсутствии ограничений по ёмкости:
\begin{equation}
M = n \cdot \bigl( p_1 + (1-p_1) p_2 + (1-p_1)(1-p_2) p_3 \bigr),
\end{equation}
$n$~--- число студентов, $p_1 = 0{,}6$, $p_2 = 0{,}25$, $p_3 = 0{,}1$. Для $n = 90$: $M \approx 65$.

\section{Решение вычислительных задач}

Генерация заявок: интервал между появлениями~--- случайная величина с равномерным распределением на отрезке $[48, 96]$ минут (среднее 72). Лимит~--- 90 заявок. Тем самым поток растягивается на первые дни недели, имитируя неравномерное прибытие студентов.

Логика потоков. Основной поток: заявки из генератора поступают в очередь, получают приоритет по $P_1$ (первая попытка~--- приоритет 1), ждут свободного канала, проходят обслуживание. Исход определяет дальнейший путь: сдача~--- выход, неудача~--- задержка 1440 минут, увеличение $P_1$, возврат в очередь. Вторичный поток~--- это повторные заявки, у них $P_1 = 2$ или $P_1 = 3$, приоритет ниже, они пропускают вперёд тех, кто сдаёт впервые. Третий элемент~--- служебный таймер: одна заявка, созданная в момент 0, выдерживает 10080 минут и по истечении времени завершает симуляцию.

\section{Анализ модели}

По результатам работы программы успешно сдали зачёт 57 из 90 студентов, то есть около 63\%. Фактически проведено 119 экзаменационных попыток из максимально возможных 120~--- ресурс зачётной недели использован почти полностью. Загрузка каналов составила 85\%, что указывает на высокую интенсивность потока.

16 человек получили слот на экзамен уже после окончания зачётной недели и были отсечены как опоздавшие. 62 студента не сдали с первой или второй попытки и к моменту завершения симуляции находились в ожидании следующего дня. Из них 5 заявок ещё ждали своей очереди на пересдачу, а 12 участников исчерпали все три попытки и покинули систему без зачёта. Таким образом, при выбранной дисциплине очереди основная часть слотов достаётся первичным сдающим, что в совокупности с ограничением по времени приводит к указанному итогу.

\refstepcounter{lstlisting}\label{lst:output}
Листинг~\ref{lst:output}. Результат работы кода.

\begin{lstlisting}
              GPSS World Simulation Report - Untitled Model 1.9.1

                   Monday, March 09, 2026 21:38:35  

           START TIME           END TIME  BLOCKS  FACILITIES  STORAGES
                0.000          10080.000    27        0          1

              NAME                       VALUE  
          AT2                            11.000
          AT3                            13.000
          ATTEMP                          3.000
          FAIL                           17.000
          GIVEUP                         22.000
          LATE                           23.000
          NXTDAY                         19.000
          PASS                        10007.000
          STOR                        10000.000
          SUCC                           14.000

 LABEL              LOC  BLOCK TYPE     ENTRY COUNT CURRENT COUNT RETRY
                    1    GENERATE            90             0       0
                    2    ASSIGN              90             0       0
ATTEMP              3    PRIORITY           135             0       0
                    4    QUEUE              135             0       0
                    5    ENTER              135             0       0
                    6    TEST               135             0       0
                    7    DEPART             119             0       0
                    8    ADVANCE            119             0       0
                    9    TEST               119             0       0
                   10    TRANSFER            78             0       0
AT2                11    TEST                41             0       0
                   12    TRANSFER            29             0       0
AT3                13    TRANSFER            12             0       0
SUCC               14    SAVEVALUE           57             0       0
                   15    LEAVE               57             0       0
                   16    TERMINATE           57             0       0
FAIL               17    LEAVE               62             0       0
                   18    TEST                62             0       0
NXTDAY             19    ADVANCE             50             5       0
                   20    ASSIGN              45             0       0
                   21    TRANSFER            45             0       0
GIVEUP             22    TERMINATE            12             0       0
LATE               23    LEAVE               16             0       0
                   24    TERMINATE            16             0       0

QUEUE              MAX CONT. ENTRY ENTRY(0) AVE.CONT. AVE.TIME   AVE.(-0) RETRY
 1                  22   16    135     53     8.985    670.851   1104.449   0

STORAGE            CAP. REM. MIN. MAX.  ENTRIES AVL.  AVE.C. UTIL. RETRY DELAY
 STOR               20   20   0    20      135   1   17.000  0.850    0    0

SAVEVALUE               RETRY       VALUE
 PASS                     0         57.000
\end{lstlisting}

\section{Практическая реализация}

\refstepcounter{lstlisting}\label{lst:gpss}
Листинг~\ref{lst:gpss}. Исходный код модели.

\begin{lstlisting}
SIMULATE

STOR    STORAGE 20

        GENERATE 72,24,,90
        ASSIGN 1,1

ATTEMP  PRIORITY P1
        QUEUE 1
        ENTER STOR,1
        TEST L AC1,8640,LATE
        DEPART 1
        ADVANCE 1440
        TEST E P1,1,AT2
        TRANSFER 0.6,FAIL,SUCC

AT2     TEST E P1,2,AT3
        TRANSFER 0.25,FAIL,SUCC

AT3     TRANSFER 0.1,FAIL,SUCC

SUCC    SAVEVALUE PASS+,1
        LEAVE STOR,1
        TERMINATE 1

FAIL    LEAVE STOR,1
        TEST L P1,3,GIVEUP
NXTDAY  ADVANCE 1440
        ASSIGN 1+,1
        TRANSFER ,ATTEMP

GIVEUP  TERMINATE 1

LATE    LEAVE STOR,1
        TERMINATE 1

        GENERATE ,,0,1
        ADVANCE 10080
        TERMINATE 1000

        RMULT 10000
        START 90
\end{lstlisting}

\section{Вывод}

В ходе работы была построена и исследована система массового обслуживания, моделирующая сдачу зачёта студентами первого курса. Выбрана очередь с приоритетом для тех, кто сдаёт впервые. Модель реализована в среде GPSS, проведён вычислительный эксперимент при 90 участниках. Получена оценка числа успешно сдавших~--- 57 человек, это ниже теоретического ожидания около 65 при неограниченной ёмкости, что объясняется лимитом 120 попыток за неделю и временным ограничением зачётного периода.
\end{document}